\documentclass[10pt,a4paper]{amsart}
\usepackage{amsmath,amssymb,amsthm,graphicx}
% For commutative diagrams we use Xy-pic
\usepackage[all]{xy}

% We use 2cell for 2-commutative diagrams.
\xyoption{2cell}
% Theorem environments.
%
\theoremstyle{plain}
\newtheorem{theorem}[subsection]{Theorem}
\newtheorem{proposition}[subsection]{Proposition}
\newtheorem{lemma}[subsection]{Lemma}

\theoremstyle{definition}
\newtheorem{definition}[subsection]{Definition}
\newtheorem{example}[subsection]{Example}
\newtheorem{exercise}[subsection]{Exercise}
\newtheorem{situation}[subsection]{Situation}

\theoremstyle{remark}
\newtheorem{remark}[subsection]{Remark}
\newtheorem{remarks}[subsection]{Remarks}

\numberwithin{equation}{subsection}

% Macros
%
\def\lim{\mathop{\mathrm{lim}}\nolimits}
\def\colim{\mathop{\mathrm{colim}}\nolimits}
\def\Spec{\mathop{\mathrm{Spec}}}
\def\Hom{\mathop{\mathrm{Hom}}\nolimits}
\def\Ext{\mathop{\mathrm{Ext}}\nolimits}
\def\SheafHom{\mathop{\mathcal{H}\!\mathit{om}}\nolimits}
\def\SheafExt{\mathop{\mathcal{E}\!\mathit{xt}}\nolimits}
\def\Sch{\mathit{Sch}}
\def\Mor{\mathop{\mathrm{Mor}}\nolimits}
\def\Ob{\mathop{\mathrm{Ob}}\nolimits}
\def\Sh{\mathop{\mathit{Sh}}\nolimits}
\def\NL{\mathop{N\!L}\nolimits}
\def\CH{\mathop{\mathrm{CH}}\nolimits}
\def\proetale{{pro\text{-}\acute{e}tale}}
\def\etale{{\acute{e}tale}}
\def\QCoh{\mathit{QCoh}}
\def\Ker{\mathop{\mathrm{Ker}}}
\def\Im{\mathop{\mathrm{Im}}}
\def\Coker{\mathop{\mathrm{Coker}}}
\def\Coim{\mathop{\mathrm{Coim}}}

% Boxtimes
%
\DeclareMathSymbol{\boxtimes}{\mathbin}{AMSa}{"02}

%
% Macros for moduli stacks/spaces
%
\def\QCohstack{\mathcal{QC}\!\mathit{oh}}
\def\Cohstack{\mathcal{C}\!\mathit{oh}}
\def\Spacesstack{\mathcal{S}\!\mathit{paces}}
\def\Quotfunctor{\mathrm{Quot}}
\def\Hilbfunctor{\mathrm{Hilb}}
\def\Curvesstack{\mathcal{C}\!\mathit{urves}}
\def\Polarizedstack{\mathcal{P}\!\mathit{olarized}}
\def\Complexesstack{\mathcal{C}\!\mathit{omplexes}}
% \Pic is the operator that assigns to X its picard group, usage \Pic(X)
% \Picardstack_{X/B} denotes the Picard stack of X over B
% \Picardfunctor_{X/B} denotes the Picard functor of X over B
\def\Pic{\mathop{\mathrm{Pic}}\nolimits}
\def\Picardstack{\mathcal{P}\!\mathit{ic}}
\def\Picardfunctor{\mathrm{Pic}}
\def\Deformationcategory{\mathcal{D}\!\mathit{ef}}
\usepackage{cleveref}
\crefname{originallemma}{original lemma}{original lemmas}
\Crefname{originallemma}{Original lemma}{Original lemmas}
\makeatletter
\newlabel{definition-bivariant-class}{{0B76}{source}{}{sourcepointer-1}{}}
\newlabel{definition-bivariant-class@cref}{{[originallemma][0][]0B76}{[1][0][]source}{}{}{}}
\makeatother
\crefname{originallemma}{original lemma}{original lemmas}
\Crefname{originallemma}{Original lemma}{Original lemmas}
\makeatletter
\newlabel{definition-envelope}{{0GU5}{source}{}{sourcepointer-2}{}}
\newlabel{definition-envelope@cref}{{[originallemma][0][]0GU5}{[1][0][]source}{}{}{}}
\makeatother
\crefname{originallemma}{original lemma}{original lemmas}
\Crefname{originallemma}{Original lemma}{Original lemmas}
\makeatletter
\newlabel{lemma-base-change-envelope}{{0GU7}{source}{}{sourcepointer-3}{}}
\newlabel{lemma-base-change-envelope@cref}{{[originallemma][0][]0GU7}{[1][0][]source}{}{}{}}
\makeatother
\crefname{originallemma}{original lemma}{original lemmas}
\Crefname{originallemma}{Original lemma}{Original lemmas}
\makeatletter
\newlabel{lemma-envelope}{{0GU8}{source}{}{sourcepointer-4}{}}
\newlabel{lemma-envelope@cref}{{[originallemma][0][]0GU8}{[1][0][]source}{}{}{}}
\makeatother
\crefname{originallemma}{original lemma}{original lemmas}
\Crefname{originallemma}{Original lemma}{Original lemmas}
\makeatletter
\newlabel{remark-restriction-bivariant}{{0F9Z}{source}{}{sourcepointer-5}{}}
\newlabel{remark-restriction-bivariant@cref}{{[originallemma][0][]0F9Z}{[1][0][]source}{}{}{}}
\makeatother
\crefname{originallemma}{original lemma}{original lemmas}
\Crefname{originallemma}{Original lemma}{Original lemmas}
\makeatletter
\newlabel{situation-setup}{{02QL}{source}{}{sourcepointer-6}{}}
\newlabel{situation-setup@cref}{{[originallemma][0][]02QL}{[1][0][]source}{}{}{}}
\makeatother
\title{Descent of a bivariant class along an envelope from the equality of its two restrictions}
\author{The Stacks Project Authors}
\date{stacks-project:chow.tex:local-snapshot-20261012; GNU FDL 1.2 or later}
\begin{document}
\maketitle
\noindent\textit{Editorial note.} All delivered mathematical text is the author's own, copied from the pinned Stacks Project chapter file chow.tex; no statement or proof step was rewritten, reordered or removed.\par\medskip
\noindent\textit{Editorial note.} Preamble provenance: the pinned chapter file opens with \textbackslash{}input\{preamble\} and carries no preamble text of its own, so the project's single Stacks-wide preamble staged at sources/stacks-chow/preamble.tex is excerpted verbatim as two spans, lines 14-18 (Xy-pic setup) and lines 154-217 (the native theorem declarations and mathematical macros). Nothing was added, removed or redefined; the wrapper class and the amsmath/amssymb/amsthm/graphicx base packages are supplied by the assembler.\par\medskip
\noindent\textit{Editorial note.} Same-chapter pointers are located at their \textbackslash{}label lines inside the pinned chapter file; cross-chapter pointers are located in the project's own tag list staged in the snapshot (sources/stacks-chow/tags.txt), at the verbatim lines carrying their canonical tags.\par\medskip
\noindent\textit{Editorial note.} The author's closing sentence at source line 6821, 'The other conditions are proved in the same manner and we omit the detailed arguments.', is retained verbatim at the end of the counted proof. Condition (1) of the bivariant class definition is carried out in full inside the span (lines 6784-6820); only the two remaining conditions of the same routine argument are left to the reader, which is the routine-verification residue that LANE\_BRIEF section 5 accepts.\par\medskip
\noindent\textit{Allowed standard prerequisites (editorial).}\begin{itemize}
\item Right exactness of the Chow sequence of an envelope. Same chapter chow.tex, label lemma-envelope, invoked at proof line 6737. Supplies that every beta in CH\_k(Y) lifts to beta' in CH\_k(Y') with g\_* beta' = beta, which is what makes the defining formula c cap beta = h\_*(c' cap beta') possible; the rest of the argument is in the counted span.
\item Restriction of a bivariant class along a cartesian square. Same chapter chow.tex, label remark-restriction-bivariant, invoked at proof lines 6732-6735. Identifies res\_1(c') and res\_2(c') with the restrictions for the two cartesian squares of the displayed diagram, which is the hypothesis consumed at proof line 6736.
\item The axioms of a bivariant class and base change of envelopes. Same chapter chow.tex, labels definition-bivariant-class (proof lines 6762 and 6782) and lemma-base-change-envelope (proof line 6766); definition-envelope (proof line 6710). Fix the three conditions that the descended operation must satisfy and guarantee that the base change Y'\_1 -> Y\_1 of an envelope is again an envelope, so the same lift can be performed after any locally finite type base change.
\end{itemize}\medskip
\ifdefined\hypertarget\hypertarget{sourcepointer-1}{}\fi
\noindent\textit{Source reference.} Original lemma 0B76: sources/stacks-chow/pinned-original-source.tex, source lines 6087--6087. Navigation pointer for \textbackslash{}ref\{definition-bivariant-class\} used inside the counted span (line 6763); the argument that consumes it is retained in the span.\par\medskip
\ifdefined\hypertarget\hypertarget{sourcepointer-2}{}\fi
\noindent\textit{Source reference.} Original lemma 0GU5: sources/stacks-chow/pinned-original-source.tex, source lines 3553--3553. Navigation pointer for \textbackslash{}ref\{definition-envelope\} used inside the counted span (line 6710); the argument that consumes it is retained in the span.\par\medskip
\ifdefined\hypertarget\hypertarget{sourcepointer-3}{}\fi
\noindent\textit{Source reference.} Original lemma 0GU7: sources/stacks-chow/pinned-original-source.tex, source lines 3580--3580. Navigation pointer for \textbackslash{}ref\{lemma-base-change-envelope\} used inside the counted span (line 6766); the argument that consumes it is retained in the span.\par\medskip
\ifdefined\hypertarget\hypertarget{sourcepointer-4}{}\fi
\noindent\textit{Source reference.} Original lemma 0GU8: sources/stacks-chow/pinned-original-source.tex, source lines 3593--3593. Navigation pointer for \textbackslash{}ref\{lemma-envelope\} used inside the counted span (line 6737); the argument that consumes it is retained in the span.\par\medskip
\ifdefined\hypertarget\hypertarget{sourcepointer-5}{}\fi
\noindent\textit{Source reference.} Original lemma 0F9Z: sources/stacks-chow/pinned-original-source.tex, source lines 6184--6184. Navigation pointer for \textbackslash{}ref\{remark-restriction-bivariant\} used inside the counted span (line 6733); the argument that consumes it is retained in the span.\par\medskip
\ifdefined\hypertarget\hypertarget{sourcepointer-6}{}\fi
\noindent\textit{Source reference.} Original lemma 02QL: sources/stacks-chow/pinned-original-source.tex, source lines 1372--1372. Navigation pointer for \textbackslash{}ref\{situation-setup\} used inside the counted span (line 6708); the argument that consumes it is retained in the span.\par\medskip
\begin{lemma}
\label{lemma-envelope-bivariant}
Let $(S, \delta)$ be as in Situation \ref{situation-setup}.
Let $f : X \to Y$ be a morphism of schemes locally of finite type over $S$.
Let $g : Y' \to Y$ be an envelope (Definition \ref{definition-envelope})
and denote $X' = Y' \times_Y X$. Let $p \in \mathbf{Z}$ and let
$c' \in A^p(X' \to Y')$. If the two restrictions
$$
res_1(c') = res_2(c') \in A^p(X' \times_X X' \to Y' \times_Y Y')
$$
are equal (see proof), then there exists a unique $c \in A^p(X \to Y)$
whose restriction $res(c) = c'$ in $A^p(X' \to Y')$.
\end{lemma}

\begin{proof}
We have a commutative diagram
$$
\xymatrix{
X' \times_X X' \ar[d]^{f''} \ar@<1ex>[r]^-a \ar@<-1ex>[r]_-b &
X' \ar[d]^{f'} \ar[r]_h &
X \ar[d]^f \\
Y' \times_Y Y' \ar@<1ex>[r]^-p \ar@<-1ex>[r]_-q &
Y' \ar[r]^g &
Y
}
$$
The element $res_1(c')$ is the restriction (see
Remark \ref{remark-restriction-bivariant}) of $c'$ for the cartesian square
with morphisms $a, f', p, f''$ and the element $res_2(c')$ is the restriction
of $c'$ for the cartesian square with morphisms $b, f', q, f''$.
Assume $res_1(c') = res_2(c')$ and let $\beta \in \CH_k(Y)$.
By Lemma \ref{lemma-envelope} we can find a $\beta' \in \CH_k(Y')$
with $g_*\beta' = \beta$. Then we set
$$
c \cap \beta = h_*(c' \cap \beta')
$$
To see that this is independent of the choice of
$\beta'$ it suffices to show that
$h_*(c' \cap (p_*\gamma - q_*\gamma))$ is zero
for $\gamma \in \CH_k(Y' \times_Y Y')$.
Since $c'$ is a bivariant class we have
$$
h_*(c' \cap (p_*\gamma - q_*\gamma)) =
h_*(a_*(c' \cap \gamma) - b_*(c' \cap \gamma)) = 0
$$
the last equality since $h_* \circ a_* = h_* \circ b_*$
as $h \circ a = h \circ b$.

\medskip\noindent
Observe that our choice for $c \cap \beta$ is forced
by the requirement that $res(c) = c'$ and the compatibility
of bivariant classes with proper pushforward.

\medskip\noindent
Of course, in order to define the bivariant class $c$ we need
to construct maps $c \cap -: \CH_k(Y_1) \to \CH_{k + p}(Y_1 \times_Y X)$
for any morphism $Y_1 \to Y$ locally of finite type satisfying the
conditions listed in Definition \ref{definition-bivariant-class}.
Denote $Y'_1 = Y' \times_Y Y_1$, $X_1 = X \times_Y Y_1$.
The morphism $Y'_1 \to Y_1$ is an envelope by
Lemma \ref{lemma-base-change-envelope}. Hence we can use the base
changed diagram
$$
\xymatrix{
X'_1 \times_{X_1} X'_1 \ar[d]^{f''_1} \ar@<1ex>[r]^-{a_1} \ar@<-1ex>[r]_-{b_1} &
X'_1 \ar[d]^{f'_1} \ar[r]_{h_1} &
X_1 \ar[d]^{f_1} \\
Y'_1 \times_{Y_1} Y'_1 \ar@<1ex>[r]^-{p_1} \ar@<-1ex>[r]_-{q_1} &
Y'_1 \ar[r]^{g_1} &
Y_1
}
$$
and the same arguments to get a well defined map
$c \cap - : \CH_k(Y_1) \to \CH_{k + p}(X_1)$ as before.

\medskip\noindent
Next, we have to check conditions (1), (2), and (3) of
Definition \ref{definition-bivariant-class} for $c$.
For example, suppose that $t : Y_2 \to Y_1$ is a proper morphism
of schemes locally of finite type over $Y$. Denote as above
the base changes of the first diagram to $Y_1$, resp.\ $Y_2$,
by subscripts ${}_1$, resp.\ ${}_2$. Denote $t' : Y'_2 \to Y'_1$,
$s : X_2 \to X_1$, and $s' : X'_2 \to X'_1$ the base changes of $t$
to $Y'$, $X$, and $X'$. We have to show that
$$
s_*(c \cap \beta_2) = c \cap t_*\beta_2
$$
for $\beta_2 \in \CH_k(Y_2)$. Choose $\beta'_2 \in \CH_k(Y'_2)$
with $g_{2, *}\beta'_2 = \beta_2$. Since $c'$ is a bivariant class
and the diagrams
$$
\vcenter{
\xymatrix{
X'_2 \ar[d]_{s'} \ar[r]_{h_2} & X_2 \ar[d]^s \\
X'_1 \ar[r]^{h_1} & X_1
}
}
\quad\text{and}\quad
\vcenter{
\xymatrix{
X'_2 \ar[d]_{s'} \ar[r]_{f'_2} & Y'_2 \ar[d]^{t'} \\
X'_2 \ar[r]^{f'_1} & Y'_1
}
}
$$
are cartesian we have
$$
s_*(c \cap \beta_2) =
s_*(h_{2, *}(c' \cap \beta'_2)) =
h_{1, *}s'_*(c' \cap \beta'_2) =
h_{1, *}(c' \cap (t'_*\beta'_2))
$$
and the final expression computes $c \cap t_*\beta_2$ by construction:
$t'_*\beta'_2 \in \CH_k(Y'_1)$ is a class whose image by $g_{1, *}$ is
$t_*\beta_2$. This proves condition (1).
The other conditions are proved in the same manner
and we omit the detailed arguments.
\end{proof}

\end{document}
